\section*{Lisad} \label{lisad} \addcontentsline{toc}{section}{Lisad}

\subsection*{I. Eksperimendis kasutatud viip} \label{lisa:promptid} \addcontentsline{toc}{subsection}{I. Eksperimendis kasutatud viip}

Käesoleva uurimuse eksperimendis kasutati kõigi nelja mudel-tööriista kombinatsiooni puhul üht ja sama lühikest eestikeelset viipa ehk ühe kasutajaviibaga (\textit{single-prompt}) lähenemist. Viipa ei muudetud ega iteratiivselt täiendatud. Täisviip oli järgmine:

\begin{quote}
\textit{``Sa oled kogenud akadeemiline kirjutaja. Loo täielik bakalaureusetöö LaTeX-vormingus (lisasin töökausta vajaliku malli), mis vastab Tartu Ülikooli arvutiteaduse instituudi nõuetele. Teema peab olema informaatikaga seotud ja vali ise. Autor on Uku Renek Kronbergs ja juhendaja Alo Peets (MSc).''}
\end{quote}

Viiba sisulised elemendid:
\begin{itemize}
    \item \textbf{Mudeli roll} -- ``kogenud akadeemiline kirjutaja'' kui üldine stiiliviide.
    \item \textbf{Ülesanne} -- täieliku bakalaureusetöö loomine.
    \item \textbf{Vorm} -- LaTeX-vorming koos viitega töökausta paigutatud UniTartuCS mallile.
    \item \textbf{Nõue} -- vastavus Tartu Ülikooli arvutiteaduse instituudi nõuetele.
    \item \textbf{Teema} -- informaatikaga seotud, valitakse mudeli poolt iseseisvalt.
    \item \textbf{Tiitellehe andmed} -- autor Uku Renek Kronbergs, juhendaja Alo Peets (MSc).
\end{itemize}

Viibasse \emph{ei} lisatud:
\begin{itemize}
    \item TÜ ATI ega LOTE lõputööde juhendite teksti;
    \item varasemate bakalaureusetööde näiteid;
    \item detailseid juhiseid peatükkide struktuuri, mahu või sisu kohta;
    \item viitamisstiili juhiseid;
    \item keelelisi korrektuurireegleid.
\end{itemize}

Sellise minimaalse viiba eesmärk oli hinnata, kuidas mudelid interpreteerivad võimalikult lühikest juhist ning millistele eelteadmistele nad toetuvad eestikeelse akadeemilise lõputöö loomisel.

\subsection*{III. Tervikliku lõputöö genereerimise näited} \label{lisa:terviktoo} \addcontentsline{toc}{subsection}{III. Tervikliku lõputöö genereerimise näited}

Järgnevalt on esitatud väljavõtted eksperimendis genereeritud terviklikest lõputöödest. Iga mudeli väljundist on toodud sissejuhatuse algus, et illustreerida erinevusi keeles, stiilis ja teemale vastavuses.

\subsubsection*{III.1 Claude Sonnet 4.6 (Cowork) -- sissejuhatuse väljavõte (40 lk)}

\begin{quote}
\textit{``Kaasaegne tarkvaraarendus on olemuselt meeskonnatöö. Isegi suurima kogemusega programmeerijad teevad vigu, ning keerukad süsteemid sisaldavad sageli turvaauke, jõudlusprobleeme või loogilisi vigu, mida üks inimene ei suuda leida.''}

\textit{``Käesoleva bakalaureusetöö peamine eesmärk on hinnata suurkeelemudelite sobivust automaatse koodiülevaatuse tööriistadena ning võrrelda eri mudeleid nende tulemuslikkuse alusel. Töö käigus viiakse läbi empiiriline katse, milles kolme LLM-i --- GPT-4, Claude 3 Opus ja Code Llama 70B --- rakendatakse 120 koodifragmendi analüüsimiseks.''}
\end{quote}

\textit{Hinnang: Töö järgib Tartu Ülikooli struktuurinõudeid ja on eestikeelne, akadeemilises stiilis. Sisaldab 23 viidet, millest 5 (22\%) olid hallutsineeritud autoriloendiga -- pealkirjad ja arXiv-identifikaatorid olid valdavalt korrektsed, kuid mudel ``täitis'' autorinimekirju välja mõeldud nimedega (kõige nähtavamad näited: \textit{Towards Automating Code Review Activities} ICSE 2021 ning \textit{Impact of Code Language Models on Automated Program Repair}, vt peatükk~\ref{subsec:sonnet-toolo}). Sarnaselt Gemini 3.1 Pro tööle on ka selle kombinatsiooni esitatud katsed (120 koodifragmendi testkogum, F1-mõõdikud GPT-4 / Claude 3 Opus / Code Llama 70B kohta) tegelikkuses kontrollimata empiirilised väited. Tiitellehel kasutab mudel autorinime kärbitud kujul (``Uku Renek''), juhendaja nimena kohatäidet (``Prof. Vanem Juhendaja, PhD'') ning aastaks ``2025''. Ristkontrollis tuvastati 37 keele- ja kirjaviga. Peamine puudus: viidete autorite hallutsineerimine, kontrollimata empiirilised väited ning kohati ülemäära keerukas lauseehitus.}

\subsubsection*{III.2 GPT-5.4 (vestlusliides) -- väljundi väljavõte (8 lk, teostatavusjuhend)}

\begin{quote}
\textit{``A complete, submission-ready bachelor's thesis written end-to-end on your behalf is not something I can produce. That would materially enable academic misconduct (submitting work largely authored by a third party). What I can provide---aligned with the University of Tartu's own principles for using generative AI in coursework---is a deep-research `thesis kit': a rigorously sourced background review, a defensible experimental design, an evaluation rubric, and a UniTartuCS-template-aligned LaTeX project skeleton with chapter-by-chapter writing prompts and clearly marked placeholders that you must complete, verify, and take responsibility for.''}
\end{quote}

\textit{Hinnang: GPT-5.4 keeldus ülesannet täitmast ja valis väljundkeeleks inglise keele. Pakkus selle asemel teostatavusjuhendit, mis sisaldas mudelivaliku soovitusi, eksperimendi disaini raamistikku ja LaTeX malli juhiseid. Viited on verifitseeritud ja sisukad. Akadeemilise aususe teadlikkus on silmapaistvalt tugev. Sama mudel sama viibaga Codex-laienduse kaudu Visual Studio Code'is genereeris aga tervikliku 47-leheküljelise eestikeelse töö -- vt III.3.}

\subsubsection*{III.3 GPT-5.4 (Codex VS Code'is) -- sissejuhatuse väljavõte (47 lk)}

\begin{quote}
\textit{``Suurte keelemudelite kiire areng on muutnud loomulikus keeles suhtlevad tarkvarasüsteemid tavapäraseks osaks nii õppetöös, ettevõtluses kui ka avalikes teenustes. [\dots] Käesolev töö keskendub eestikeelsele dokumentide põhisele küsimus-vastus süsteemile.''}

\textit{``Töö eesmärk on analüüsida RAG-lahenduste teaduskirjandust ning selle põhjal kavandada eestikeelse dokumentide põhise küsimus-vastus süsteemi arhitektuur, mis oleks läbipaistev, ajakohastatav ja hinnatav.''}
\end{quote}

\textit{Hinnang: Sama mudel ja sama viip kui vestlusliideses, kuid agentse töövoo kaudu andis täiesti erineva tulemuse. Keelevalik eesti keel, struktuur vastab bakalaureusetöö vormile (8 nummerdatud peatükki, 49 alapeatükki, 7 tabelit, 2 joonist), ortograafiline täpsus silmapaistvalt kõrge (3 kirjaviga kogu töö peale). Kõik 22 unikaalset viidet olid täielikult kontrollitavad ja korrektsete metaandmetega -- ainus genereeritud töö, kus ühtegi hallutsineeritud viidet ei tuvastatud. Tähelepanuväärne kõrvalleid: tiitellehel, litsentsis ja allkirjas kasutas mudel viibas etteantud nimede asemel kohatäiteid (\textit{Näidisüliõpilane}, \textit{Juhendaja nimi, PhD}), vältides reaalse tudengi nime alla allakirjastatud akadeemilise teksti otsest tootmist.}

\subsubsection*{III.4 Gemini 3.1 Pro (Antigravity) -- sissejuhatuse väljavõte (39 lk)}

\begin{quote}
\textit{``Tarkvaraarenduse valdkond on viimaste aastate jooksul teinud läbi märkimisväärse arengutõuke tänu erinevate tehisintellekti vormide laialdasele levikule ja kättesaadavusele. Üha enam integreeritakse generatiivset tehisintellekti (GenAI) igapäevastesse tööriistadesse ja arenduskeskkondadesse (IDE).''}

\textit{``Käesoleva bakalaureusetöö eesmärk on just selle andmelünga osaline täitmine kohalikus kontekstis. Töö keskendub Eesti informaatikaliõpilastele ja püüab selgitada välja, kuidas täpselt mõjutab lokaalse tehisintellekti viipeassistendi kasutamine algaja programmeerija produktiivsust esimesel õppeaastal.''}
\end{quote}

\textit{Hinnang: Gemini 3.1 Pro + Antigravity kombinatsioon genereeris vormiliselt korrektse töö generatiivse TI mõjust tarkvaraarendusele -- väidetavalt 2026.~aasta märtsis Tartu Ülikooli Delta õppehoone arvutiklassis läbi viidud kvaasi-eksperimendi näitel, milles osales 30 informaatika esmakursuslast (16~meest, 14~naist; 18--22~aastat). Kuigi teema on informaatikaga seotud ja seega viibale vastav, on kogu selles kirjeldatud empiiriline osa -- 30 tudengiga kvaasi-eksperiment, testikeskkond Monaco Editor'i ja Docker'i liivakastiga, kohaliku LLaMA-3-8B-Instruct mudeli kasutamine TÜ HPC klastris ning kõik statistilised tulemused -- fabritseeritud: reaalseid tudengeid, katseid ega andmeid ei eksisteeri. Eestikeelne tekst on ladus, kuid sisaldab palju ilmseid kirjavigu (24 ristkontrollitud leidu). Seda juhtumit käsitletakse töös akadeemilise väärkasutuse kõrge riskiga juhtumina (vt peatükk~\ref{subsec:gemini-toolo}).}

\subsubsection*{III.5 Claude Opus 4.6 (Cowork) -- sissejuhatuse väljavõte (61 lk)}

\begin{quote}
\textit{``Tänapäeva digimaailmas on igapäevaelu lahutamatult seotud veebiteenustega. Inimesed kasutavad regulaarselt kümneid, sageli isegi sadu erinevaid platvorme --- alates e-postist ja sotsiaalmeediast kuni pangateenuste ja tervishoiuportaalideni.''}

\textit{``Käesoleva bakalaureusetöö eesmärk on kavandada ja arendada turvaline paroolihalduri veebirakendus, mis põhineb null-teadmise arhitektuuril. Rakendus kasutab kaasaegseid krüptograafilisi algoritme --- AES-256-GCM sümmeetrilist krüpteerimist ja Argon2id võtmetuletust --- ning on loodud avatud lähtekoodina, et tagada läbipaistvus ja auditeeritavus.''}
\end{quote}

\textit{Hinnang: Opus 4.6 + Cowork kombinatsioon genereeris kõige mahukama ja tehniliselt sügavaima töö. 61-leheküljeline lõputöö sisaldas terviklikku tehnilist arhitektuuri ja koodi\-näiteid; lisaks \emph{väitis} töö, et rakendus läbis 149 automatiseeritud testi, OWASP Top 10 analüüsi ja SUS kasutatavusuuringu (skoor 77,0). Neid teste, auditit ega kasutajauuringut ei teostanud käesoleva uurimuse autor iseseisvalt -- tegemist on kombinatsiooni enda genereeritud empiiriliste väidetega, mis tuleb inimese poolt kontrollida samamoodi nagu viited. Eestikeelne akadeemiline stiil oli läbivalt suurepärane (0 esimese isiku asesõna terve töö peale); ristkontrollis tuvastati 15 keele- ja kirjaviga. Struktuur vastas TÜ nõuetele. 35-st viitest 33 olid täielikult korrektsed; 2 viites (6\%) oli väike metaandmete viga (URL-i kuju ning pealkirja/aasta ebatäpsus). Peamine eelis teiste kombinatsioonide ees: tehniline detailsus, \LaTeX{} koodi otsene genereerimine failidesse ja võime hallata terviklikku projektistruktuuri.}

\subsubsection*{III.6 Genereeritud tööde võrdlev kokkuvõte}

\begin{table}[ht]
\centering
\caption{Tervikliku lõputöö genereerimise koondvõrdlus}
\label{tab:terviktoo-vordlus}
\begin{tblr}{
    width=\linewidth,
    colspec={X[1.5,l] X[1,c] X[1,c] X[1,c] X[1,c]},
    rows={font=\small},
    colsep=3pt,
    row{1}={font=\bfseries\small},
    hlines,
    vlines,
    row{even}={bg=colorCellHighlight},
}
Kriteerium & Gemini 3.1 Pro + Antigravity & GPT-5.4 + Codex & Claude Sonnet 4.6 + Cowork & Claude Opus 4.6 + Cowork \\
Keel & Eesti & Eesti & Eesti & Eesti \\
Maht & 39 lk & 47 lk & 40 lk & 61 lk \\
Valitud teema & GenAI ja progr. & RAG QA-süsteem & LLM koodiülevaatuses & Paroolihaldur \\
Struktuur & Lõputöö & Lõputöö & Lõputöö & Lõputöö \\
Keele- ja kirjavigu (kogu töö peale) & 24 & 3 & 37 & 15 \\
Unikaalseid viiteid & 9 & 22 & 23 & 35 \\
Hallutsineeritud viiteid & 1 (11\%) & 0 (0\%) & 5 (22\%) & 2 (6\%) \\
Esimese isiku asesõnu & 1 & 0 & 0 & 0 \\
Praktiline kasutus aluspõhjana & Madal (fabritseeritud empiirika) & Kõrge & Kõrge & Väga kõrge (kontrollimata empiiriliste väidetega) \\
\end{tblr}
\end{table}
